\section{Related Work}
\label{sec:related-work}

Reconstructing a collision from vehicle recordings connects accident
investigation with cooperative perception and multisensor estimation.
We consider how these approaches recover motion and relate vehicle
identities and recordings in time to support forensic interpretation.

\begin{table*}[!t]
\caption{Closest forensic reconstruction systems. Entries summarize the
reported procedures and demonstrations.}
\label{tab:closest-reconstruction}
\centering
\footnotesize
\renewcommand{\arraystretch}{1.15}
\begin{tabular*}{\textwidth}{@{\extracolsep{\fill}}p{0.17\textwidth}p{0.25\textwidth}p{0.23\textwidth}p{0.28\textwidth}@{}}
\hline
\textbf{Work} & \textbf{Recorded evidence} & \textbf{Temporal alignment} & \textbf{Reconstruction and interpretation} \\
\hline
Beck et al.~\cite{beck2023pipeline} &
Camera, LiDAR, radar; vehicle logs &
Synchronized inputs; recovery outside scope &
Image/depth tracking of a selected target \\
Matsumura et al.~\cite{matsumura2024precrash} &
EDR + dashcam from one or both vehicles &
Chronological EDR/video integration &
Trajectories of known crash participants \\
Zhang et al.~\cite{zhang2025reverse} &
Driving logs + dashcam/map &
Vehicle-log timestamps &
Motion from logged relative distances \\
Yuan et al.~\cite{yuan2026multisource} &
EDR + GPS + roadside video &
External time reference and correction &
Kalman filtering; vehicle-state and scene-parameter fusion \\
\hline
\textbf{Ours} &
Raw radar + vehicle logs  &
Clock recovery from shared contacts or observed motion &
Cross-observer association; Kalman filtering and RTS smoothing;
semantic events with recorder provenance \\
\hline
\end{tabular*}
\end{table*}

\subsection{Accident Reconstruction from Recorded Evidence}

Dynamics models provide a basis for reconstructing vehicle motion and
contacts in simulation tools such as PC-Crash~\cite{steffan1996pccrash},
with EDR measurements offering independent checks on the resulting
estimates~\cite{gabler2004edr}. Recorded perception and exchanged vehicle
states provide complementary routes to motion reconstruction.
DeepCrashTest uses trajectories extracted from monocular dashcam videos
to create virtual crash tests~\cite{bashetty2020deepcrashtest}, while Young
et al. reconstruct host and neighbouring trajectories from an EDR storing
states received through cooperative collision-warning
messages~\cite{young2009reconstruction}.

\subsection{Cooperative Perception and Tracking}

Combining communicated object tracks in high-level cooperative
perception~\cite{rauch2012car2x} requires correspondences between vehicles'
object lists~\cite{rauch2013association} as well as spatial and temporal
alignment~\cite{allig2019alignment}. Differences in observation time are
addressed through asynchronous automotive track fusion by Aeberhard
et al.~\cite{aeberhard2012asynchronous}, time-compensated late fusion in
DAIR-V2X~\cite{yu2022dairv2x}, and cooperative detection under asynchronous
sharing and pose errors in V2X-ViT~\cite{xu2022v2xvit}.

Uncertainty is addressed by Chiu et al., who learn detection uncertainty
within a multisensor Kalman framework for cooperative tracking~\cite{chiu2024dmstrack},
while Castelino et al. compare track-fusion strategies using Collective
Perception Messages in asynchronous scenarios, assuming reliable source
timestamps~\cite{castelino2026fusion}. Together, these works establish
asynchronous fusion and cross-vehicle association as central tasks in
cooperative perception.

\subsection{Association, Fusion, and Temporal Registration}

Correlated estimation errors affect both track association and
fusion~\cite{barshalom1981correlation,barshalom1986process}; covariance
intersection provides an established way to combine estimates when
their cross-correlations are unknown~\cite{julier1997intersection}.
For automotive radar, extended-object tracking also models the multiple
reflection centres on one vehicle~\cite{hammarstrand2012extended}, while
consistent cooperative localization addresses uncertainty in relative
pose and correlation~\cite{hery2021localization}. These principles inform
our association tests, which ask whether one rigid body can explain the
surfaces seen by different observers while accounting for uncertainty
and competing identities. Later evidence can then establish earlier
identities while preserving the original local observations and event times.

Common scene dynamics~\cite{caspi2002alignment} and passive sensor timing
estimation~\cite{olson2010synchronization} likewise provide precedents for
temporal registration. In the present method, shared collision reports
constrain constant clock offsets after checks against recorded geometry
and available radar motion. When this evidence cannot align a recorder,
its offset is estimated by comparing radar observations of a recording
vehicle with that vehicle's own logged motion and footprint. Temporal
registration is thus coupled to the physical evidence used for vehicle
association.

\subsection{Forensic Timelines and Interpretation}

The NTSB's Tempe investigation demonstrates the forensic value of relating
an automated vehicle's perception history to its subsequent
response~\cite{ntsb2019tempe}. Response models relate action timing to
the evolving situation and driver expectations~\cite{engstrom2024response},
while counterfactual studies such as Scanlon et al.'s reconstructed fatal
crashes assess alternative outcomes through additional
simulation~\cite{scanlon2021fatal}.
These applications motivate preserving the connection between
reconstructed interactions and recorded responses in a forensic chronology.

\subsection{Positioning of This Work}
\label{subsec:positioning}

Table~\ref{tab:closest-reconstruction} compares how the closest forensic
systems turn recorded evidence into aligned motion and reconstruction
outputs. Image and depth tracking of a selected target allow Beck
et al.~\cite{beck2023pipeline} to reconstruct uninstrumented surrounding
vehicles, with clock synchronization outside their scope. Other systems
integrate vehicle records with scene information:
Matsumura et al.~\cite{matsumura2024precrash} combine EDR trajectories and
dashcam frames from one or both crash participants in a chronological
road map, while Zhang et al.~\cite{zhang2025reverse} reproduce traffic
scenes in OpenSCENARIO from timestamped driving logs and relative
distances. Explicit temporal alignment is central to Yuan
et al.~\cite{yuan2026multisource}, who align EDR, GPS and roadside video
to an external time reference and use Kalman filtering to handle noise
and timing deviations.

In our setting, both clock offsets and physical vehicle correspondences
must be recovered across independent radar recordings. Shared contacts
and observed motion constrain relative time, and body geometry connects
the different surfaces seen by different observers. Motion is estimated
through Kalman filtering and RTS smoothing, while sustained
cross-observer evidence can also resolve earlier identity fragments.
Together, these steps support reconstruction of vehicles with and
without their own logs, including road users absent from the impact
reports, without requiring camera-based correspondence or communication
during the incident.

The semantic chronology makes this reconstruction inspectable across
recorders by relating manoeuvres and conflicts to recorded responses
and contacts, while retaining local observations, event times and the
provenance of later identity decisions. The contribution therefore lies
in integrating temporal and physical correspondence recovery with
retrospective motion estimation and a traceable semantic account.
The evaluation examines these aspects separately through collision
identification, physical association, geometry and local event fidelity
under controlled sensor degradation.
